存储账单压过了日志本身的价值，于是有人提议只留每天的计数与求和。这个提议成不成立，取决于两个问题：删掉细节以后，还能不能算出原来要算的结论；就算什么都不删，这批数据又最多能把参数说到多准。本单元把这两个问题接成一条链——先问压缩能否无损，再问精度的上限在哪，最后把上限写成一个可以拿到会上用的数字。

第一步是分清哪些数据细节对参数是多余的。给定统计量 \(T\) 后样本的条件分布若不再依赖 \(\theta\)，则原始数据对该参数不再有话可说，这就是充分性；因子分解定理把这个判断变成对联合密度的一次代数观察，指数族则解释了为什么固定长度的摘要在实践中如此常见，从而使流式聚合与分片合并成为可能。这是\hyperref[sec:11-Mathematical-Statistics/03-Sufficient-Statistics-and-Information/01]{``充分统计量保留了哪些参数信息''}。

第二步是把``这份数据对参数说得有多死''量化。对数似然对参数的导数称为 score，它在真值处平均为零，方差就是 Fisher 信息；同一个量还等于对数似然的平均负曲率，于是似然曲面越尖，参数越容易辨认，越平则数据说不清。向量参数下信息矩阵的退化方向，直接对应实验设计没有覆盖到的方向。这是\hyperref[sec:11-Mathematical-Statistics/03-Sufficient-Statistics-and-Information/02]{``Score 与 Fisher 信息衡量局部可辨认性''}。

第三步是把曲率换算成方差的硬下限。任何无偏估计量的方差都不低于 Fisher 信息的倒数，推导只用了一次协方差恒等式与一次 Cauchy--Schwarz，正则条件与无偏性在其中各出场一次。这条界的用处是分诊：估计量还在下界上方，换方法有收益；已经贴住下界，就只能加数据或改模型设定。这是\hyperref[sec:11-Mathematical-Statistics/03-Sufficient-Statistics-and-Information/03]{``Cramér--Rao 界给无偏估计设下限''}。

有三处容易串味的地方值得预先标出。Fisher 信息不是 Shannon 熵换个名字，它是局部的、依赖参数化的量，度量的是参数微小移动后分布有多容易被区分；充分性与低维压缩是两件事，原始数据本身永远充分，却一个字节都没省；Cramér--Rao 界只约束无偏估计，有偏估计可以走到界的下面，正则化改善均方误差的空间正来自这个缺口。三篇里反复出现的正则条件也不是装饰，支持集随参数移动的模型会让其中大部分结论整块失效。
